\section{Open contractions with native edge orientations}
\label{sec:peps_oriented_open_contractions}

Let $\Gamma=(V,E)$ be a finite simple graph with ordered vertices.
Write an ordered edge as $e=(t,h)$, with $t<h$, and let
$o:E\to\{0,1\}$ specify its orientation: $o_e=0$ means $t\to h$
and $o_e=1$ means $h\to t$. The ordering of vertices is used only
for coordinates. The physical pair on an edge is always ordered as
$(x_h,x_t)$, even when the native arrow is reversed.
For $R\subseteq V$, denote its internal, crossing, and exterior edges by
$E_R$, $\partial E_R$, and $E_{R^c}$.
For a region family $\mathcal R$, write $\mathcal P_{\mathcal R}(A)$
for the physical vectors whose every slice on each $R\in\mathcal R$
belongs to the regional contraction range of $A$.
These conventions extend the bond transformations of
\cite[Section~7]{Schuch2010PEPS} to the rightward and downward arrows
of a torus.

\subsection{Averaging tensors and unrestricted regional boundaries}

\begin{definition}[Native tails]%
    \label{def:peps_oriented_native_tail}
    \lean{TNLean.PEPS.graphNativeTail,
        TNLean.PEPS.graphNativeTail_edgeLeftIncident,
        TNLean.PEPS.graphNativeTail_edgeRightIncident}
    \leanok
    An incidence $(v,e)$ is a native tail when $v=t$ and $o_e=0$,
    or when $v=h$ and $o_e=1$.
    Thus the ordered first incidence is a native tail exactly when
    $o_e=0$, and the ordered second incidence exactly when $o_e=1$.
\end{definition}

\begin{definition}[Oriented averaging tensor]%
    \label{def:peps_oriented_averaging_tensor}
    \lean{TNLean.PEPS.graphOrientedAveragingSite,
        TNLean.PEPS.graphOrientedAveragingSite_false,
        TNLean.PEPS.numberedGraphOrientedAveragingTensor,
        TNLean.PEPS.numberedGraphOrientedAveragingTensor_false}
    \leanok
    \uses{def:peps_oriented_native_tail}
    Let $G$ be finite, $U:G\to\mathrm{GL}(\C^X)$ a representation,
    and $W$ a matrix on $\C^X$. For virtual and physical incidence labels
    $\eta,\sigma:\operatorname{Inc}(v)\to X$, set
    \begin{align}
        A^{U,W,o}_v(\eta,\sigma)
        =\frac1{|G|}\sum_{g\in G}
        \prod_{e\ni v}
        \begin{cases}
            (WU_{g^{-1}})_{\eta_e,\sigma_e}, & v\text{ is the native tail}, \\
            (U_gW)_{\sigma_e,\eta_e},        & v\text{ is the native head}.
        \end{cases}
        \label{eq:peps_oriented_site}
    \end{align}
    A bijection $e_v:X^{\operatorname{Inc}(v)}\to\{0,\ldots,p-1\}$
    numbers the physical indices. Its coefficient at $s$ is
    $A^{U,W,o}_v(\eta,e_v^{-1}(s))$. Fix any numbering of $X$ on the
    virtual bonds. If every $o_e=0$, this is the ordered
    dressed averaging tensor. No symmetry of $W$ is required.
    The normalization is the group-average normalization of
    \cite{gap:rmp_peps_quantum_double_g_isometry}.
\end{definition}

\begin{definition}[Oriented open bond coefficients]%
    \label{def:peps_oriented_open_coefficients}
    \lean{TNLean.PEPS.graphOrientedOpenBoundaryFactor,
        TNLean.PEPS.graphOrientedOpenInternalFactor,
        TNLean.PEPS.graphOrientedOpenBondCoordinates,
        TNLean.PEPS.graphOrientedOpenBondSpace,
        TNLean.PEPS.graphOrientedOpenBoundaryFactor_false,
        TNLean.PEPS.graphOrientedOpenInternalFactor_false,
        TNLean.PEPS.graphOrientedOpenBondCoordinates_false,
        TNLean.PEPS.graphOrientedOpenBondSpace_false}
    \leanok
    \uses{def:peps_oriented_averaging_tensor}
    Fix $q:R\to G$. On an internal edge $e=(t,h)$ define
    \begin{align}
        I_e(q;x_h,x_t) & =
        \begin{cases}
            (W^2U_{q_hq_t^{-1}})_{x_h,x_t}, & o_e=0, \\
            (W^2U_{q_tq_h^{-1}})_{x_t,x_h}, & o_e=1.
        \end{cases}
        \label{eq:peps_oriented_internal}
    \end{align}
    For a crossing edge, let $v$ be its retained vertex, $\theta$ its
    virtual boundary label, and $x$ its retained physical label. Put
    $B_e(q;\theta,x)=(WU_{q_v^{-1}})_{\theta,x}$ when $v$ is the
    native tail, and $B_e(q;\theta,x)=(U_{q_v}W)_{x,\theta}$ otherwise.
    For arbitrary $\theta\in X^{\partial E_R}$, define
    \begin{align}
        C_R^{U,W,o}(\theta)(x,z)
         & =|G|^{-|R|}\sum_{q:R\to G}
        \prod_{e\in\partial E_R}B_e(q;\theta_e,x_e)
        \prod_{e\in E_R}I_e(q;z_e),
        \label{eq:peps_oriented_open}                 \\
        \mathcal S_R(U,W,o)
         & =\operatorname{span}\{C_R^{U,W,o}(\theta):
        \theta\in X^{\partial E_R}\}.
        \label{eq:peps_oriented_open_span}
    \end{align}
    Here $z_e=(x_h,x_t)$ and $x$ lists the retained crossing endpoints.
    Setting $o=0$ recovers each ordered boundary factor, internal factor,
    open vector, and open span.
\end{definition}

\begin{theorem}[The complete regional contraction range]%
    \label{thm:peps_oriented_open_range}
    \lean{TNLean.PEPS.graphOrientedOpenBondCoordinates_eq_sum_prod,
        TNLean.PEPS.graphOpenRegionNetwork_orientedAveragingSite_coherent,
        TNLean.PEPS.graphOpenBondRegrouping_orientedAveragingSite,
        TNLean.PEPS.openRegionWeight_numberedGraphOrientedAveragingTensor,
        TNLean.PEPS.regionBondPhysicalEquiv_openRegionWeight_oriented,
        TNLean.PEPS.map_regionGroundSpace_eq_graphOrientedOpenBondSpace,
        TNLean.PEPS.regionBondPhysicalEquiv_mem_graphOrientedOpenBondSpace_iff}
    \leanok
    \uses{def:peps_oriented_open_coefficients,def:peps_region_ground_space}
    Assume $WU_g=U_gW$ for every $g$. Let $\mathcal C_R$ be the
    open contraction of the tensors~(\ref{eq:peps_oriented_site}),
    and let $J_R$ regroup its physical incidences into crossing labels
    and ordered internal pairs. Then
    \begin{align}
        J_R\mathcal C_R(\theta)          & =C_R^{U,W,o}(\theta),
        \label{eq:peps_oriented_regroup}                         \\
        J_R\operatorname{im}\mathcal C_R & =\mathcal S_R(U,W,o).
        \label{eq:peps_oriented_range}
    \end{align}
    These identities hold before or after numbering physical and virtual
    labels. In particular, $\psi\in\operatorname{im}\mathcal C_R$
    exactly when $J_R\psi\in\mathcal S_R(U,W,o)$.
    The boundary coefficients may be arbitrary linear combinations of
    all $\theta$, including entangled boundary vectors.
\end{theorem}
\begin{proof}\leanok
    Expand the normalized group average separately at every vertex,
    obtaining the factor $|G|^{-|R|}$ and the sum over $q:R\to G$.
    Regroup the incidence product by internal and crossing edges.
    Summation over one internal virtual index gives either
    $(U_{q_h}W)(WU_{q_t^{-1}})$ or
    $(U_{q_t}W)(WU_{q_h^{-1}})$ with the physical indices interchanged.
    Commutation with $W$ gives~(\ref{eq:peps_oriented_internal}).
    Each crossing edge retains its original $\theta_e$, giving
    (\ref{eq:peps_oriented_open}); commuting the scalar products also
    gives the internal-first form of the same sum.
    The bijections of incidence and numbered labels preserve every finite
    sum. Finally, the regional range is the span of all its boundary
    columns, proving~(\ref{eq:peps_oriented_range}).
\end{proof}

\begin{definition}[Oriented incident representation]%
    \label{def:peps_oriented_incident_representation}
    \lean{TNLean.PEPS.graphOrientedIncidentMatrix,
        TNLean.PEPS.graphOrientedIncidentMatrix_one,
        TNLean.PEPS.graphOrientedIncidentMatrix_mul,
        TNLean.PEPS.graphOrientedIncidentRepresentation,
        TNLean.PEPS.graphOrientedIncidentRepresentation_apply,
        TNLean.PEPS.graphOrientedNumberedIncidentRepresentation,
        TNLean.PEPS.graphOrientedNumberedIncidentRepresentation_apply}
    \leanok
    \uses{def:peps_oriented_native_tail}
    On $\C^{X^{\operatorname{Inc}(v)}}$, define
    \begin{align}
        \rho_v^o(g)_{\sigma,\eta}
        =\prod_{e\ni v}
        \begin{cases}
            U_{g^{-1}}(\eta_e,\sigma_e), & v\text{ is the native tail}, \\
            U_g(\sigma_e,\eta_e),        & v\text{ is the native head}.
        \end{cases}
        \label{eq:peps_oriented_incident}
    \end{align}
    Matrix multiplication gives $\rho_v^o(1)=1$ and
    $\rho_v^o(gh)=\rho_v^o(g)\rho_v^o(h)$. At a tail this follows
    from $(U_{(gh)^{-1}})^T=(U_{g^{-1}})^T(U_{h^{-1}})^T$.
    Thus these matrices act as a representation by multiplication on
    coefficient vectors. If $E_v$ changes numbered virtual coefficients
    to incidence coefficients, the numbered representation is
    $E_v^{-1}\rho_v^o(g)E_v$.
\end{definition}

\begin{theorem}[$G$-injectivity of the oriented average]%
    \label{thm:peps_oriented_average_g_injective}
    \lean{TNLean.PEPS.regularSiteMap_graphOrientedAveragingSite,
        TNLean.PEPS.isGInjective_graphOrientedAveragingSite,
        TNLean.PEPS.localTensorMap_numberedGraphOrientedAveragingTensor_one,
        TNLean.PEPS.isGInjective_numberedGraphOrientedAveragingTensor_one}
    \leanok
    \uses{def:peps_oriented_incident_representation,
        def:peps_oriented_averaging_tensor,def:peps_g_injective}
    The local coefficient map of $A^{U,1,o}_v$ is
    $\Pi_v=|G|^{-1}\sum_g\rho_v^o(g)$, and is $G$-injective.
    With physical numbering $P_v$ and virtual numbering $E_v$, its
    local map is $P_v\Pi_v E_v$ and is $G$-injective for
    $E_v^{-1}\rho_v^o E_v$.
    This assertion needs neither unitarity nor semi-regularity of $U$.
\end{theorem}
\begin{proof}\leanok
    Compare~(\ref{eq:peps_oriented_site}) at $W=1$ with
    (\ref{eq:peps_oriented_incident}) coefficient by coefficient.
    The group average satisfies $\Pi_v\rho_v^o(g)=\Pi_v$ and fixes
    every invariant vector. Hence its restriction to the invariant
    subspace is injective. The invertible coordinate changes $P_v,E_v$
    preserve both assertions.
\end{proof}

\subsection{Internal support and regional blocks}

\begin{theorem}[Oriented internal factors determine parent support]%
    \label{thm:peps_oriented_parent_support}
    \lean{TNLean.PEPS.graphOrientedOpenInternalFactor_mem_range_of_factors,
        TNLean.PEPS.graphRegionInternalBondSlice_orientedAveragingSite,
        TNLean.PEPS.graphRegionInternalBondSlice_mem_range_of_oriented_factors,
        TNLean.PEPS.regionInternalBondSlice_mem_range_of_oriented_factors,
        TNLean.PEPS.regionParentGroundSpace_mem_product_range_of_oriented_factors}
    \leanok
    \uses{def:peps_oriented_open_coefficients,
        def:peps_region_parent_hamiltonian}
    Let $F:\C^Z\to\C^{X\times X}$. Assume both coefficient vectors
    of $W^2U_g$ and $(W^2U_g)^T$ belong to $\operatorname{im}F$
    for every $g$. Then each $I_e(q)$ belongs to $\operatorname{im}F$.
    Assume also $WU_g=U_gW$. Fixing the physical crossing labels $x$
    in an open vector leaves
    \begin{align}
        |G|^{-|R|}\sum_q
        \left(\prod_{e\in\partial E_R}B_e(q;\theta_e,x_e)\right)
        \bigotimes_{e\in E_R}I_e(q)
        \ \in\operatorname{im}F^{\otimes E_R}.
        \label{eq:peps_oriented_internal_support}
    \end{align}
    This holds for every vector in the regional range after taking its
    internal physical slice. If a region family $\mathcal R$ contains
    both endpoints of every edge in some member, every vector in its
    parent space $\mathcal P_{\mathcal R}(A^{U,W,o})$ lies, after
    regrouping, in $\operatorname{im}F^{\otimes E}$.
\end{theorem}
\begin{proof}\leanok
    \uses{thm:peps_oriented_open_range,thm:peps_parent_support_product_range}
    Each orientation selects one of the two assumed matrix orders.
    Choose a preimage under $F$ for every $I_e(q)$ and take their
    product, with the scalar coefficient in
    (\ref{eq:peps_oriented_internal_support}). Its image under
    $F^{\otimes E_R}$ is the displayed sum.
    Span closure proves the assertion for every regional vector.
    For a global parent vector, choose a region containing the endpoints
    of a prescribed edge. Every one-edge slice then belongs to
    $\operatorname{im}F$. The product-range criterion gives the
    full $E$-fold range.
\end{proof}

\begin{definition}[Products of edge-dependent physical maps]%
    \label{def:peps_oriented_bond_family}
    \lean{TNLean.PEPS.graphNumberedBondFamilyMatrix,
        TNLean.PEPS.graphBoundaryBondFamilyMatrix,
        TNLean.PEPS.graphExteriorBondFamilyCoefficient,
        TNLean.PEPS.graphNumberedBondFamilyMatrix_eq_submatrix}
    \leanok
    For $F_e:\C^{X\times X}\to\C^{Y\times Y}$, define
    $M(F)=\bigotimes_{e\in E}F_e$, expressed in the numbered site
    coordinates through the incidence-to-bond bijections.
    Thus $M(F)_{\tau,\sigma}=\prod_eF_e((\tau_h,\tau_t),(\sigma_h,\sigma_t))$
    in bond coordinates. Fix exterior endpoint labels $\gamma_Y,\gamma_X$.
    On a crossing edge put
    \begin{align}
        F_{e,R}^{\gamma_Y,\gamma_X}(y,x)
        =\begin{cases}
             F_e((\gamma_{Y,e},y),(\gamma_{X,e},x)), & t\in R, \\
             F_e((y,\gamma_{Y,e}),(x,\gamma_{X,e})), & h\in R.
         \end{cases}
        \label{eq:peps_oriented_crossing_block}
    \end{align}
    For fixed exterior physical configurations $\tau,\sigma$, let
    $c_R(F;\tau,\sigma)$ be the product of their matrix coefficients
    on all $e\in E_{R^c}$.
\end{definition}

\begin{theorem}[Regional blocks of an edge-family product]%
    \label{thm:peps_oriented_bond_family_blocks}
    \lean{TNLean.PEPS.regionOperatorBlock_graphNumberedBondFamilyMatrix}
    \leanok
    \uses{def:peps_oriented_bond_family}
    The block obtained by fixing exterior output $\tau$ and input
    $\sigma$ satisfies
    \begin{align}
        M(F)_{R;\tau,\sigma}
        =c_R(F;\tau,\sigma)\,
        J_{Y,R}^{-1}\left(
        \bigotimes_{e\in\partial E_R}F_{e,R}^{\gamma_Y,\gamma_X}
        \otimes\bigotimes_{e\in E_R}F_e\right)J_{X,R},
        \label{eq:peps_oriented_family_block}
    \end{align}
    where $\gamma_Y,\gamma_X$ are the crossing endpoint labels determined
    by $\tau,\sigma$, and $J_{X,R},J_{Y,R}$ regroup regional coordinates.
\end{theorem}
\begin{proof}\leanok
    Partition the defining product for $M(F)$ into exterior, crossing,
    and internal edges. Exterior edges give $c_R$; fixing the outside
    endpoint of a crossing edge gives
    (\ref{eq:peps_oriented_crossing_block}). The internal factors are
    unchanged. The coordinate bijections identify their product with
    (\ref{eq:peps_oriented_family_block}).
\end{proof}

\begin{theorem}[Composition and adjoints of edge-family products]%
    \label{thm:peps_oriented_bond_family_algebra}
    \lean{TNLean.PEPS.graphNumberedBondFamilyMatrix_conjTranspose,
        TNLean.PEPS.graphNumberedBondFamilyMatrix_mul,
        TNLean.PEPS.graphNumberedBondFamilyMatrix_one,
        TNLean.PEPS.graphNumberedBondFamilyMatrix_mul_eq_one}
    \leanok
    \uses{def:peps_oriented_bond_family}
    With compatible finite endpoint alphabets,
    $M(F)^\dagger=M(F^\dagger)$,
    $M(F)M(L)=M((F_eL_e)_e)$, and $M((1)_e)=1$.
    In particular, $F_eL_e=1$ on every edge implies $M(F)M(L)=1$.
\end{theorem}
\begin{proof}\leanok
    Conjugation distributes over each coefficient product.
    For composition, the sum over the intermediate site configuration
    is, under the bond-coordinate bijection, a sum over independent
    intermediate edge labels. It factors into the products $F_eL_e$.
    The product of edgewise Kronecker deltas is the delta of the
    complete configuration, proving the identity assertion.
\end{proof}
